\section{Hartogs extension for boxes}\label{sec:hartogs}

Our first extension step is local in the source but uses the compact
K\"ahler target.  A general meromorphic extension can acquire exceptional
fibers.  For a box, the two separate tangencies exclude these fibers, even
when the box has deficient rank.

\begin{proposition}[Hartogs extension]\label{prop:hartogs}
Let $(s,t)$ be the standard coordinates of $\C^2$, and let $(w,v)$ be any
invertible complex-affine coordinate system.  For real numbers
$0<e<a$ and $0<c<b$, put
\[
 H=\bigl\{|w|<e,\ |v|<b\bigr\}
   \ \cup\ \bigl\{|w|<a,\ c<|v|<b\bigr\},
 \qquad
 Q=\bigl\{|w|<a,\ |v|<b\bigr\}.
\]
Every holomorphic map $f:H\to X$ satisfying
\[
 \dd f(\partial_s)\in E_1,
 \qquad
 \dd f(\partial_t)\in E_2
\]
extends uniquely to a holomorphic map $Q\to X$ with the same tangencies.
Thus the affine coordinates defining the Hartogs figure need not be the
coordinates defining the box.
\end{proposition}

\begin{proof}
We first obtain a meromorphic extension.  We then rule out curves in its
graph fibers by comparing two types of level sets: a level isolated at a
putative exceptional curve, and nearby levels containing a parameter
direction.  Finally we show that the remaining finite fibers are singletons.

\emph{Meromorphic extension and tangency.}
The envelope of holomorphy of $H$ is $Q$.  Indeed, expand a scalar
holomorphic function on the shell in a Laurent series in $v$.  The
coefficients of the negative powers are holomorphic functions of $w$;
they vanish for $|w|<e$, and hence for $|w|<a$.  Cauchy estimates on any
intermediate circle $c<|v|<b$ give normal convergence of the resulting
power series on compact subsets of $Q$.  This is the usual Hartogs
extension calculation; since the bidisk $Q$ is a domain of holomorphy, it
also identifies the envelope.

Ivashkovich's extension theorem states that a meromorphic map from a
domain in a Stein manifold to a compact K\"ahler manifold extends
meromorphically to the envelope of holomorphy
\cite[Theorem~2.2]{Ivashkovich2008}; see also \cite{Ivashkovich1992}.
Here $H$ is a domain in the Stein manifold $\C^2$, and $f$ is holomorphic,
so the theorem gives a meromorphic extension to $Q$.

Let $\Gamma\subset Q\times X$ be its graph, with its reduced analytic
structure, and let $\rho:\Gamma\to Q$ be the first projection.  The graph
is an irreducible analytic surface, $\rho$ is proper, and there is a proper
analytic subset $A\subset Q$, disjoint from $H$, such that over
$U=Q\setminus A$ the graph is the ordinary graph of a holomorphic map
$F:U\to X$.  This ordinary graph is dense in $\Gamma$.  The set $U$ is
connected: the complement of a proper complex analytic subset of a
connected complex manifold is connected.  Consequently, the identity
theorem applied to the two holomorphic bundle-valued expressions gives
\begin{equation}\label{eq:meromorphic-box-tangency}
 P_2\dd F(\partial_s)=0,
 \qquad
 P_1\dd F(\partial_t)=0
 \quad\hbox{on }U.
\end{equation}
Here $P_i:T_X\to E_i$ is the holomorphic projection determined by the
specified splitting.

\emph{No positive-dimensional fiber.}
We use the local irreducible decomposition, dimension theorem, and
regular-locus properties of complex analytic sets; see
\cite[Chapter~II, Section~4]{DemaillyCADG}.
We record first a point that is important when $\Gamma$ is singular or
locally reducible.  On every local irreducible surface branch of $\Gamma$,
both source coordinates $s$ and $t$ are nonconstant.  To see this, choose
a relatively open part of that branch away from the other local branches.
It meets the dense ordinary graph.  Near such a point, $\rho$ is an
isomorphism onto an open subset of $Q$, and $s,t$ are source coordinates.
Neither coordinate can therefore be constant on the branch.  It follows
that, for every $s_0$, the analytic set
\begin{equation}\label{eq:graph-coordinate-level}
 \Gamma\cap\{s=s_0\}
\end{equation}
has local dimension at most one; the same holds with $t$ in place of $s$.

Suppose a fiber over $(s_0,t_0)$ has positive dimension.  Since it is
contained in \eqref{eq:graph-coordinate-level}, it has a curve component
$C$.  Take a small piece of $C$ in a product chart
$z=(z_1,z_2)$ of $X$.  Some scalar component $h$ of $z_1$ or $z_2$ is
nonconstant on this piece: otherwise the fiber curve, which lies in
$\{(s_0,t_0)\}\times X$, would be a point.

First suppose $h$ belongs to $z_1$.  The curve piece is a component of
the reduced set \eqref{eq:graph-coordinate-level}, since that set has
dimension at most one.  Choose a general point $q$ of the
curve piece such that the reduced set \eqref{eq:graph-coordinate-level}
is a smooth curve near $q$ and $h$ is nonconstant on its germ.  This is
possible because the one-dimensional analytic set has locally finitely
many curve components, distinct components meet discretely, and the
singular points of a reduced curve are discrete.  In particular, $q$ is
an isolated point of the level set
\[
 (s,h)^{-1}\bigl(s_0,h(q)\bigr)\subset\Gamma.
\]
This assertion concerns the reduced coordinate level, and does not
require $\Gamma$ itself to be smooth along $C$.

Work in ambient complex coordinates on $Q\times X$ around $q$, and choose
a sufficiently small closed Euclidean ball $\overline B$ centered at $q$.
Its closure lies in this coordinate neighborhood, and the indicated
level set misses $\Gamma\cap\partial B$.  The image of this boundary
intersection under the holomorphic map
\[
 \Psi=(s,h)
\]
is compact and does not contain $\Psi(q)$.  Choose points $q_j$ of the
ordinary graph tending to $q$.  For all sufficiently large $j$, the level
value $\Psi(q_j)$ also misses that compact image.  Therefore
\[
 K_j=\Gamma\cap\overline B\cap\Psi^{-1}\bigl(\Psi(q_j)\bigr)
\]
is a compact analytic subset of the open ball $B$: it has no points on
$\partial B$, so there is no boundary truncation of the analytic level.
On the ordinary graph near $q_j$, however,
\eqref{eq:meromorphic-box-tangency} says that $h$ is independent of $t$.
Varying $t$ with $s$ fixed gives a one-dimensional germ in $K_j$ through
$q_j$.  This contradicts the fact that a compact analytic subset of an
open set in Euclidean complex space has no positive-dimensional
component.  One way to see the latter fact is to maximize the strictly
plurisubharmonic squared coordinate norm on such a component and restrict
it to a nonconstant analytic curve germ through a maximum point.

If $h$ belongs to $z_2$, the identical argument uses $\Psi=(t,h)$ and
varies $s$.  Thus every fiber of $\rho$ has dimension zero.  Properness
then makes each fiber finite.

\emph{Singleton fibers and holomorphic removal.}
Fix $p\in Q$, and write its fiber as
$\{(p,x_1),\ldots,(p,x_k)\}$.  Choose pairwise disjoint target
neighborhoods $V_1,\ldots,V_k$ with pairwise disjoint closures and
$x_i\in V_i$.  Properness implies that, after shrinking a base ball
$D$ about $p$, every point of $\rho^{-1}(D)$ has its target coordinate in
$\bigcup_iV_i$.  Otherwise a sequence of graph points over points tending
to $p$, outside this union, would have a limit in the central fiber.
The ordinary locus $D\setminus A$ is connected, so its image under $F$
lies in a single $V_i$.  Density of the ordinary graph then places every
$x_j$ in $\overline V_i$.  The disjoint closures imply $k=1$.

The proper bijection $\rho$ is a homeomorphism.  Its inverse followed by
the second projection defines a continuous extension
$\widehat F:Q\to X$.  Near each source point its image lies in a target
coordinate chart.  The coordinate functions are continuous and
holomorphic off $A$, hence holomorphic by the removable-singularity
theorem for locally bounded holomorphic functions across analytic
subsets of a complex manifold.  Thus $\widehat F$ is holomorphic.
Equation~\eqref{eq:meromorphic-box-tangency} extends by the identity
theorem, giving the required box tangencies.  Uniqueness follows from
the identity theorem on the connected bidisk $Q$.
\end{proof}
